% Worksheet W7 – Refine and evaluate (2 pages)
% Übersetzen: latexmk -xelatex ws-w7-refine-and-evaluate.tex
% Lösungsfassung: ws-w7-refine-and-evaluate-lsg.tex
% Zahlen stammen aus den Übungsdaten der Website (englisch/w-seminar-time/data/*.csv):
% survey-pilot.csv (tp3 umgepolt), coding-round1.csv, test-pilot.csv, experiment-pilot.csv, experiment-main.csv
% ============================================================
% ab-vorlage.tex – gemeinsame Vorlage der Arbeitsblätter
% englisch/w-seminar-time/arbeitsblaetter/  (W-Seminar "Time")
% Übersetzen mit:  latexmk -xelatex <datei>.tex   (oder zweimal xelatex)
% Gestaltung wie englisch/13/questions-on-the-text/arbeitsblaetter/ab-vorlage.tex,
% Antwortplatz-Makros (8-mm-Linien) wie informatik/10/java/arbeitsblaetter/ab-vorlage.tex
% Schriften: Carlito (Fließtext), Poppins (Überschriften)
% ============================================================
\documentclass[10pt,a4paper]{article}
\usepackage[a4paper,top=12mm,bottom=13mm,left=14mm,right=14mm,headsep=4mm,headheight=12pt,footskip=7mm]{geometry}
\usepackage{fontspec}
\setmainfont{Carlito}
\newfontfamily\headfont{Poppins}[UprightFont=* Medium, BoldFont=* Medium]
\newfontfamily\headlight{Poppins Light}
\usepackage{xcolor}
\definecolor{prim}{HTML}{2563EB}
\definecolor{primdark}{HTML}{1E3A8A}
\definecolor{sec}{HTML}{7C3AED}
\definecolor{dusk}{HTML}{171A3F}
\definecolor{brass}{HTML}{D9AB2F}
\definecolor{primtint}{HTML}{EFF6FF}
\definecolor{sectint}{HTML}{F5F3FF}
\definecolor{ink}{HTML}{1E293B}
\definecolor{muted}{HTML}{64748B}
\definecolor{rule}{HTML}{CBD5E1}
\definecolor{good}{HTML}{059669}
\definecolor{goodtint}{HTML}{ECFDF5}
\definecolor{warn}{HTML}{B45309}
\definecolor{warntint}{HTML}{FFFBEB}
\definecolor{bad}{HTML}{B91C1C}
\definecolor{badtint}{HTML}{FEF2F2}
% Farben der Zeit-Dimensionen (wie auf der Website)
\definecolor{dlit}{HTML}{B83280}
\definecolor{dphys}{HTML}{2563EB}
\definecolor{dcs}{HTML}{0F8F84}
\definecolor{dpsy}{HTML}{C26A06}
\definecolor{dlang}{HTML}{7C3AED}
\definecolor{dbey}{HTML}{4B5563}
\color{ink}
\usepackage[most]{tcolorbox}
\usepackage{tikz}
\usetikzlibrary{positioning,calc,arrows.meta,shapes.misc}
\usepackage{enumitem}
\setlist{nosep,leftmargin=1.1em,topsep=1pt}
\setlist[itemize]{label=\textcolor{prim}{\small\textbullet}}
\usepackage{tabularx,array,booktabs,colortbl}
\usepackage{multicol}
\usepackage{fancyhdr}
\usepackage{lastpage}
\usepackage{amssymb}
\usepackage{calc}
\usepackage{qrcode}
\usepackage{microtype}
\setlength{\parindent}{0pt}
\setlength{\parskip}{0pt}
\linespread{1.02}

% ---------- Kopf- und Fußzeile ----------
\newcommand{\wskopf}{W-Seminar English · Time · Year 12}
\pagestyle{fancy}
\fancyhf{}
\renewcommand{\headrulewidth}{0pt}
\fancyhead[L]{\footnotesize\headlight\textcolor{muted}{\wskopf}}
\fancyhead[R]{\ifnum\value{page}=1\footnotesize\headlight\textcolor{muted}{Name: \rule{38mm}{0.3pt}\quad Date: \rule{18mm}{0.3pt}}\fi}
\fancyfoot[L]{\scriptsize\textcolor{muted}{edu-mrh.de/englisch/w-seminar-time}}
\fancyfoot[R]{\scriptsize\textcolor{muted}{\thepage\,/\,\pageref*{LastPage}}}

% ---------- Titel ----------
% \wstitel{Kürzel}{Titel}{Untertitel}
\newcommand{\wstitel}[3]{%
  \begin{tcolorbox}[enhanced,colback=dusk,colframe=dusk,arc=2mm,boxrule=0pt,
     left=4mm,right=4mm,top=2.2mm,bottom=2.2mm,
     interior style={left color=dusk,right color=primdark}]
    \color{white}{\headlight\footnotesize #1}\\[0.3mm]
    {\headfont\LARGE #2}\\[0.6mm]
    {\small #3}
  \end{tcolorbox}\vspace{1mm}}

% ---------- Boxen ----------
% \begin{wsbox}[Optionen]{Nummer}{Überschrift} ... \end{wsbox}
\newtcolorbox{wsbox}[3][]{enhanced,breakable,colback=white,colframe=rule,boxrule=0.5pt,arc=1.5mm,
  left=2.5mm,right=2.5mm,top=1.3mm,bottom=1.3mm,
  title={\headfont\normalsize\textcolor{white}{\colorbox{prim}{\,#2\,}}\;\textcolor{ink}{#3}},
  coltitle=ink,colbacktitle=white,attach title to upper={\par\vspace{0.6mm}},before skip=1.4mm,after skip=1.4mm,#1}
\newtcolorbox{tintbox}[1][]{enhanced,colback=primtint,colframe=primtint,boxrule=0pt,arc=1.2mm,
  left=2mm,right=2mm,top=1.2mm,bottom=1.2mm,#1}
\newtcolorbox{secbox}[1][]{enhanced,colback=sectint,colframe=sectint,boxrule=0pt,arc=1.2mm,
  left=2mm,right=2mm,top=1.2mm,bottom=1.2mm,#1}
\newtcolorbox{goodbox}[1][]{enhanced,colback=goodtint,colframe=goodtint,boxrule=0pt,arc=1.2mm,
  left=2mm,right=2mm,top=1.2mm,bottom=1.2mm,#1}
\newtcolorbox{warnbox}[1][]{enhanced,colback=warntint,colframe=warntint,boxrule=0pt,arc=1.2mm,
  left=2mm,right=2mm,top=1.2mm,bottom=1.2mm,#1}
\newtcolorbox{badbox}[1][]{enhanced,colback=badtint,colframe=badtint,boxrule=0pt,arc=1.2mm,
  left=2mm,right=2mm,top=1.2mm,bottom=1.2mm,#1}

\newcommand{\kl}[1]{{\headfont\small\textcolor{prim}{#1}}}      % kleine Zwischenüberschrift
\newcommand{\klsec}[1]{{\headfont\small\textcolor{sec}{#1}}}
\newcommand{\klc}[2]{{\headfont\small\textcolor{#1}{#2}}}       % Zwischenüberschrift in Farbe
\newcommand{\plus}{\textcolor{good}{\textbf{+}}\ }
\newcommand{\minus}{\textcolor{bad}{\textbf{–}}\ }
\newcommand{\pfeil}{\textcolor{prim}{$\rightarrow$}\ }
\newcommand{\cb}{\raisebox{-0.2ex}{\textcolor{muted}{$\square$}}\ }   % Ankreuzkästchen

% ---------- Antwortplatz (8 mm, wie informatik/10/java) ----------
\newif\ifloesung
\ifdefined\mitloesung \loesungtrue \else \loesungfalse \fi
\newcommand{\lsgfarbe}{\color{sec}\itshape}
\newlength{\zeilenabstand}\setlength{\zeilenabstand}{8mm}
\tikzset{schreiblinie/.style={draw=black!45,line width=0.35pt}}
% Lücke: \luecke[Breite] / \lsg[Breite]{Text}
\newcommand{\lsg}[2][45mm]{%
  \makebox[#1][l]{\rlap{\raisebox{-0.6ex}{\rule{#1}{0.4pt}}}\ifloesung{\lsgfarbe\small\,#2}\fi}}
\newcommand{\luecke}[1][45mm]{\lsg[#1]{}}
% Lücke bis Zeilenende
\newcommand{\lsgrest}[1]{\leavevmode\unskip\hspace{1.5mm}\ifloesung\rlap{\raisebox{0.8pt}{\lsgfarbe\,#1}}\fi\leaders\hbox{\rule[-0.6ex]{0.5pt}{0.4pt}}\hfill\null\par}
\newcommand{\restzeile}{\lsgrest{}}
% nur in der Lösung sichtbar (Platz bleibt frei), z. B. in Tabellenzellen
\newcommand{\lsgtext}[1]{\ifloesung\leavevmode{\lsgfarbe\small#1}\fi}
\newcommand{\lsgoder}[2]{\ifloesung\leavevmode#1\else#2\fi}
\newcommand{\lueckentext}{\linespread{1.5}\selectfont\raggedright}
\newcommand{\schreibzeile}{\rule[-2.8mm]{0pt}{8mm}}
% Schreiblinien: \linien{n} / \lsglinien{n}{Text}; erste Linie eine volle Zeilenhöhe unter der Frage
\newcommand{\lsglinien}[2]{\par\vspace{0.5mm}\noindent\begin{tikzpicture}
  \path (0,0) (\linewidth-0.4pt,-#1\zeilenabstand-1.5mm);
  \foreach \i in {1,...,#1}{\draw[schreiblinie] (0.2pt,-\i\zeilenabstand) -- (\linewidth-0.2pt,-\i\zeilenabstand);}
  \ifloesung\node[overlay,anchor=north west,inner sep=0pt] at (1mm,-\zeilenabstand+1pt+5mm) {\parbox[t]{\linewidth-2mm}{\lsgfarbe\raggedright\setlength{\baselineskip}{\zeilenabstand}\rule{0pt}{5mm}#2}};\fi
  \end{tikzpicture}\par}
\newcommand{\linien}[1]{\lsglinien{#1}{}}

\newcommand{\website}{https://edu-mrh.de/englisch/w-seminar-time/data-refine.html}
\newcommand{\cbx}{\lsgoder{\raisebox{-0.2ex}{\textcolor{sec}{\boldmath$\boxtimes$}}\ }{\cb}}
\newcommand{\phase}[2]{\colorbox{sectint}{\headfont\scriptsize\textcolor{sec}{#1}\,\textcolor{muted}{#2}}}
\newcommand{\buchst}[1]{\fbox{\rule{0pt}{2.2ex}\makebox[4mm]{\lsgtext{#1}}}}   % Kästchen für einen Buchstaben

\begin{document}

\wstitel{Worksheet W7 · W-Seminar “Time”}{Refine and evaluate}%
{Test your instrument in a pilot, read the signals – and say what your data can and cannot show.}

% ------------------------------------------------------------------ 1
\begin{wsbox}{1}{Read the pilot: a time-pressure scale}
{\small 24 students answered six statements from 1 (strongly disagree) to 5 (strongly agree); tp3 is already reversed. \textbf{α with all six items = 0.71.}}\par\vspace{0.8mm}
\renewcommand{\arraystretch}{1.05}
{\small
\begin{tabularx}{\linewidth}{@{}l X c c c c@{}}
\arrayrulecolor{rule}\toprule
 & \textbf{statement} & \textbf{mean} & \textbf{SD} & \textbf{item–rest r} & \textbf{α without it}\\\midrule
tp1 & I often feel I don’t have enough time for my schoolwork. & 2.96 & 1.04 & 0.58 & 0.63\\
tp2 & My days feel rushed. & 3.08 & 1.02 & 0.70 & 0.59\\
tp3 & I usually have time to relax in the evening. \textcolor{muted}{(reversed)} & 3.08 & 1.21 & 0.49 & 0.67\\
tp4 & I feel stressed and tired in the morning. & 3.38 & 0.71 & $-$0.18 & 0.80\\
tp5 & Deadlines often catch me by surprise. & 3.33 & 1.27 & 0.72 & 0.57\\
tp6 & I sometimes have more to do than I can manage. & 4.58 & 0.50 & 0.44 & 0.70\\\bottomrule
\end{tabularx}}
\par\vspace{0.5mm}{\scriptsize\textcolor{muted}{tp6: 10 students ticked 4, 14 ticked 5, nobody ticked 1, 2 or 3.}}\par\vspace{0.5mm}
\begin{lueckentext}
a) I would drop or split item \lsg[22.5mm]{tp4} because\par
\lsgrest{it asks two things at once – its item–rest r is even negative}
b) Without it, α rises to \lsg[22.5mm]{0.80}.\quad c) tp6 shows a \lsg[52.5mm]{ceiling effect}.\par
Better wording for tp6: \lsgrest{Most days I have more to do than I can manage.}
\end{lueckentext}
\end{wsbox}

% ------------------------------------------------------------------ 2
\begin{wsbox}{2}{More signals \enspace{\normalfont\small\textcolor{muted}{one per data type}}}
\begin{lueckentext}
\kl{Coding:} two coders agree on 77\,\% of the lines, by chance 20\,\%.\enspace κ = (0.77 − 0.20) / (1 − 0.20) = \lsg[22.5mm]{0.71}\newline
Good enough to code the whole corpus? \lsgrest{no – aim for 0.80: add a decision rule to the codebook}
\kl{Test:} item q4 is solved by 97\,\% of the students: \lsgrest{too easy – it separates nobody}
Item q7 has a \emph{negative} discrimination. First check: \lsgrest{the key (a second correct answer?)}
\kl{Experiment:} two estimates of a 3-minute wait are 210 and 165. \lsgrest{seconds – rewrite the instruction}
The fast music was rated 3.6, the slow music 3.3 (1–5): \lsgrest{manipulation too weak – stronger contrast}
\end{lueckentext}
\end{wsbox}

% ------------------------------------------------------------------ 3
\begin{wsbox}{3}{My pilot \enspace{\normalfont\small\textcolor{muted}{plan it, run it, log every change}}}
\begin{minipage}[t]{0.8\linewidth}\vspace{0pt}
\begin{lueckentext}
\kl{Who, how many, when:}\enspace\lsgrest{5 students from Year 9, next Tuesday}
\end{lueckentext}
{\small\kl{I watch}\enspace \cbx time needed\quad \cbx questions they ask\quad \cb think-aloud\quad \cbx impossible values\par\vspace{0.5mm}
  \hspace*{12mm}\cb α / κ / item analysis\quad \cbx manipulation check\quad \cb other: \luecke[30mm]}
\end{minipage}\hfill
\begin{minipage}[t]{0.17\linewidth}\vspace{0pt}\raggedleft
\qrcode[height=15mm]{\website}
\end{minipage}
\par\vspace{1mm}
{\small\phase{Pair}{3 min} Tell your partner what is most likely to go wrong in your pilot. Note their idea:}
\begin{lueckentext}
\lsgrest{people know the song – ask afterwards and note it}
\end{lueckentext}
\vspace{0.5mm}
\renewcommand{\arraystretch}{1}
\begin{tabularx}{\linewidth}{@{}X X X@{}}
\arrayrulecolor{rule}\toprule
\textbf{what I saw} & \textbf{what I changed} & \textbf{why}\\\midrule
\rule[-4mm]{0pt}{11mm}\lsgtext{2 of 5 asked “minutes or seconds?”} & \lsgtext{instruction with example “2 min 40 s”} & \lsgtext{no more impossible values}\\\midrule
\rule[-4mm]{0pt}{11mm} & & \\\midrule
\rule[-4mm]{0pt}{11mm} & & \\\bottomrule
\end{tabularx}
\end{wsbox}

% ================================================================== page 2
\newpage
% ------------------------------------------------------------------ 4
\begin{wsbox}{4}{Which analysis answers the question?}
{\small Choose: \textbf{compare} groups · \textbf{relate} two variables · track a \textbf{change} · \textbf{count} categories.}\par\vspace{0.8mm}
\renewcommand{\arraystretch}{1}
\begin{tabularx}{\linewidth}{@{}X p{22mm} p{30mm} p{30mm}@{}}
\arrayrulecolor{rule}\toprule
\textbf{research question} & \textbf{analysis} & \textbf{chart} & \textbf{key number}\\\midrule
\schreibzeile Do evening types sleep less on school nights? & \lsgtext{relate} & \lsgtext{scatter plot} & \lsgtext{r (or ρ)}\\\midrule
\schreibzeile Does fast music make a wait feel shorter than slow music? & \lsgtext{compare} & \lsgtext{box plot} & \lsgtext{d and p}\\\midrule
\schreibzeile Has the share of “gonna” in song lyrics grown since the 1960s? & \lsgtext{count} & \lsgtext{line chart (shares)} & \lsgtext{\% per decade}\\\midrule
\schreibzeile Did a grammar lesson improve test scores more than in a control class? & \lsgtext{change} & \lsgtext{box plot of gains} & \lsgtext{mean gain, d}\\\midrule
\schreibzeile \textbf{My question:} & & & \\\bottomrule
\end{tabularx}
\par\vspace{1mm}
{\small\phase{Share}{1 min each} Explain your analysis plan to your partner: which result would surprise you – and what would you do then?}
\end{wsbox}

% ------------------------------------------------------------------ 5
\begin{wsbox}{5}{Read the output and write it up}
\begin{tintbox}
{\ttfamily\small fast: M = 2.60 min (SD 0.61), n = 28 \quad slow: M = 3.04 min (SD 0.62), n = 28\\
difference 0.45 min, d = 0.73 \quad permutation test: p = 0.009}
\end{tintbox}
\begin{lueckentext}
a) The effect is \lsg[52.5mm]{medium to large} \enspace(d ≈ 0.2 small · 0.5 medium · 0.8 large).\par
\end{lueckentext}\par\vspace{0.5mm}
b) p = 0.009 means, in your own words:
\lsglinien{2}{If the music made no difference, a gap this large would appear in fewer than 1 in 100 shuffles.}
c) Your results sentence (numbers, no interpretation yet):
\lsglinien{3}{Students with fast music estimated the wait as shorter (M = 2.60, SD = 0.61) than students with slow music (M = 3.04, SD = 0.62); d = 0.73, p = .009.}
d) One limitation:
\lsglinien{2}{One school, volunteers only – and we measured estimated, not felt, time.}
\end{wsbox}

% ------------------------------------------------------------------ 6
\begin{wsbox}{6}{Found nothing? Diagnose before you conclude}
{\small A pilot with 12 students: screen time before bed and time pressure, \textbf{r = 0.04}.\enspace
\phase{Think}{2 min} \phase{Pair}{3 min} Find three reasons why r could be close to 0 although there is a link – and decide what to do.}
\par\vspace{0.5mm}
\begin{lueckentext}
1 \lsgrest{too few people (n = 12) – r jumps around}
2 \lsgrest{unreliable measure, e.g. one vague item}
3 \lsgrest{restricted range / U-shape / subgroups cancel out – plot it}
\kl{In the pilot I would:}\enspace\lsgrest{fix the instrument (more items, wider sample), pilot again}
\kl{In the main study I would:}\enspace\lsgrest{report r, n and effect honestly; discuss the causes}
\end{lueckentext}
\end{wsbox}

\end{document}
